\scanpage{13}
\deleted{\unsure{Morphismes quasi-projectifs}}

\unsure{Épim. couvrants|Épimorphismes|épimorphismes universels} (Zar) de schémas (= foncteurs représentables)

% slanted
\margin{\textcircled{b} faisceaux sur (Sch)}

\textcircled{d} Schém. de type fini sur un autre \inserted{$S$} : [quasi-comp. et sur $S$ et loc. de type fini sur $S$]

\textcircled{a} Schémas comme foncteurs : avantages respectifs des arguments anneaux (schémas affines) et des arguments schémas (\uline{voire espaces localement annelés}).

\textcircled{e} Corollaire du th. de représentabilité : nature \unsure{localement|locale} du \unsure{pb.|Pb.}

\textcircled{f} Schéma $P(M)$, $M$ module qu.-coh sur $S$ \inserted{(pas néc. affine)} : il représente un foncteur remarquable \unsure{m.} dans la catégorie des espaces loc. annelés sur $S$ (voire des topos loc. annelés sur $S$). Sa structure essentielle : --

% slanted
\margin{c'est un th. de représentabilité}

Dépendance p.r. au changement de base ; dépendance fonctorielle en $S$ / Dépendance en $M$ (pour épi)

[Schémas \inserted{quasi-}projectifs sur $S$ ($S$ affine disons) : qui se plongent \deleted{\unsure{d}} dans un $P(M)$, \deleted{\unsure{de de de présentation}} \inserted{(on peut alors choisir $V$ de type fini)}]

\inserted{de t.f. sur $S$,}

Relation avec $V(M)$ (fibré vectoriel sur $S$ associé à $M$)

\[ V(M)^*\subset V(M)\qquad(\text{les \uline{épi}}\quad M_{k'}\longrightarrow k') \]

\[ V(M)^*\longrightarrow P(M)\qquad\text{\uline{épi de faisceaux}}\quad[\text{\uline{i.e. il y a une section loc pour Zar.}}] \]

En fait, le $\unsure{\mathrm{Gr}_m|G_m|\mathbb{G}_{m,S}}$ opère sur $V(M)^{\deleted{*}}$, laissant $V(M)^*$ stable ; il y opère librement, et

\[ P(M)\simeq V(M)^*/G_m\quad(\text{isom de faisceaux Zar}) \]

\textcircled{g} $\varprojlim$ finies dans (Sch)

\textcircled{h} Spectre d'Algèbre qu. coh. et \deleted{\unsure{schéma}} morphismes affines (les schémas affines sur $S$)

% archive sheet number: 12
